
\chapter{语义分析}


语义分析(semantic analysis)是在词法分析和句法分析的基础上, 进一步识别文本所表达的意义内涵.

如果说词法分析回答“文本里有哪些词”, 句法分析回答“这些词怎样组成句子”,
那么语义分析要回答的是:

\begin{Verbatim}[breaklines=true,breakanywhere=true,fontsize=\small]
语言符号与客观世界中的事件、实体、属性和关系如何对应
\end{Verbatim}

在 NLP 流水线中, 语义分析位于词法和句法处理之后:

\begin{quote}
\small
\textbf{语义分析在 NLP 流水线中的位置}
\par 节点: sem-raw: 原始文本; sem-lex: 词法分析; sem-syn: 句法分析; sem-sem: 语义分析; sem-understand: 文本理解与推理; sem-lex-out: 词、词性、实体; sem-syn-out: 结构树或依存树; sem-sem-out: 词义、论元、角色
\par 边: sem-raw -\textgreater{} sem-lex; sem-lex -\textgreater{} sem-syn; sem-syn -\textgreater{} sem-sem; sem-sem -\textgreater{} sem-understand; sem-lex -\textgreater{} sem-lex-out; sem-syn -\textgreater{} sem-syn-out; sem-sem -\textgreater{} sem-sem-out
\end{quote}

语义分析的内容很多, 任务也很艰难.
课件主要聚焦两类典型任务:

\begin{itemize}
\item 词义消歧(word sense disambiguation, WSD).
\item 语义角色标注(semantic role labeling, SRL).
\end{itemize}

\section{语义分析概述}


\subsection{什么是语义分析}


课件中的定义可以概括为:
对输入文本内容, 在词法分析和句法分析基础上, 建立计算模型,
让计算机识别语言符号背后所表达的意义及其与客观世界的真实联系.

这个定义包含两个层次:

\begin{itemize}
\item 表层语言符号: 词、短语、句法结构等可直接观察的形式.
\item 深层语义联系: 词义、事件、参与者、指代关系、篇章关系、隐含意图等.
\end{itemize}

语义分析比词法和句法更困难, 因为意义往往不是由单个词直接决定的,
而要结合上下文、常识、说话场景和任务目标.

\subsection{任务类型划分}


按分析对象的层次, 语义分析可分为:

\begin{table}[H]
\centering
\begingroup
\setlength{\tabcolsep}{2pt}
\small
\caption{语义分析任务的层次划分}
\begin{tabular}{|p{0.28\textwidth}|p{0.28\textwidth}|p{0.28\textwidth}|}
\hline
层面 & 典型任务 & 核心问题 \\ \hline
词语层面 & 词义消歧 & 多义词在当前上下文中取哪个义项 \\ \hline
句子层面 & 语义角色标注 & 谓词与句中成分之间是什么语义关系 \\ \hline
篇章层面 & 共指消解、话题结构分析 & 不同句子中的实体和话题如何关联 \\ \hline
\end{tabular}
\endgroup
\end{table}

按分析深度, 又可分为浅层语义分析和深层语义推理:

\begin{itemize}
\item 浅层语义分析: 使用符号表面层之下一层的隐含关系, 如词义、语义角色、局部事件结构.
\item 深层语义推理: 使用更深层的信息和内在关系, 如常识、意图、事实蕴含和语境反讽.
\end{itemize}

例如“你是一个好人。”在字面上是正面评价;
但在特定对话场景中, 可能表达委婉拒绝.
这种差异说明深层语义往往需要结合语境和常识推理.

\begin{quote}
\small
\textbf{从表层语言符号到深层语义推理}
\par 节点: sem-surface: 语言符号表面层; sem-shallow: 词义与句内关系; sem-deep: 篇章、常识与意图; sem-shallow-task: 浅层语义分析; sem-deep-task: 深层语义推理
\par 边: sem-surface -\textgreater{} sem-shallow; sem-shallow -\textgreater{} sem-deep; sem-shallow -\textgreater{} sem-shallow-task; sem-deep -\textgreater{} sem-deep-task
\end{quote}

\section{词义消歧}


\subsection{基本概念}


词义消歧是指:
确定一个多义词在给定上下文语境中的具体含义.

例如:

\begin{Verbatim}[breaklines=true,breakanywhere=true,fontsize=\small]
The fish was bought by the cook.
The fish was bought by the river.
\end{Verbatim}

第一句中的 \texttt{fish} 更自然地理解为“鱼肉/鱼类商品”,
第二句中的 \texttt{fish} 则更可能指“鱼这种生物”.
又如 \texttt{bank} 在不同上下文中可表示“银行”或“河岸”.

词义消歧需要和词性标注区分:

\begin{itemize}
\item 词性标注判断一个词在当前句子中是什么词类, 如名词、动词、形容词.
\item 词义消歧判断同一词类下或不同词类下的具体义项.
\end{itemize}

词性不同的多义词通常较容易处理, 因为词性标注已经提供了强线索.
真正困难的是词性相同、语义又相近的多义词, 因为它们往往需要更丰富的上下文和常识信息.

\subsection{多义现象类型}


课件把词的多义现象分为三类:

\begin{table}[H]
\centering
\begingroup
\setlength{\tabcolsep}{2pt}
\small
\caption{多义词的三类情况}
\begin{tabular}{|p{0.28\textwidth}|p{0.28\textwidth}|p{0.28\textwidth}|}
\hline
类型 & 例子 & 说明 \\ \hline
意义相关的多义 & \texttt{open}: “开着的”/“公开的” & 义项之间有关联, 边界较模糊 \\ \hline
意义无关的多义 & \texttt{bank}: “银行”/“河岸” & 义项来源或含义差异大 \\ \hline
词性不同的多义 & \texttt{cook}: “煮”/“厨师” & 义项常伴随词性差异 \\ \hline
\end{tabular}
\endgroup
\end{table}

上下文信息对不同任务的作用并不相同:

\begin{itemize}
\item 对词性标注, 邻近上下文往往很有效, 例如前后几个词能强烈提示当前词类.
\item 对词义消歧, 相隔较远的实词常常更有效, 因为主题词和语义场能提示当前词义.
\end{itemize}

因此, WSD 不能只看目标词附近的局部词形, 还要考虑整句甚至篇章中的主题线索.

\subsection{方法分类}


课件列出的主要方法包括:

\begin{quote}
\small
\textbf{词义消歧方法分类}
\par 节点: wsd-root: 词义消歧; wsd-rule: 基于规则; wsd-supervised: 有监督学习; wsd-dict: 基于词典; wsd-rule-sub: 选择限制; wsd-mi: 互信息; wsd-ml: 贝叶斯 / 最大熵; wsd-dict-sub: 定义 / 义类 / 双语词典
\par 边: wsd-root -\textgreater{} wsd-rule; wsd-root -\textgreater{} wsd-supervised; wsd-root -\textgreater{} wsd-dict; wsd-rule -\textgreater{} wsd-rule-sub; wsd-supervised -\textgreater{} wsd-mi; wsd-supervised -\textgreater{} wsd-ml; wsd-dict -\textgreater{} wsd-dict-sub
\end{quote}

这些方法的本质都是利用上下文信息, 判断当前语境更支持哪一个义项.

\subsection{基于规则的方法}


基于规则的 WSD 根据语言学研究结果, 人工抽取词与词连用的限制规则,
再用这些规则对词义进行判定.

一个典型规则是选择限制(selectional restriction).
例如英文 \texttt{have} 若取“吃”的含义, 通常要求:

\begin{itemize}
\item 主语具有动物性或生命性.
\item 宾语具有可食用性.
\end{itemize}

课件中还给出动词“跑”的例子:

\begin{Verbatim}[breaklines=true,breakanywhere=true,fontsize=\small]
快速行走: 主语 +动物 -> 加高分
行驶:     主语 +车辆 -> 加高分
\end{Verbatim}

对于句子:

\begin{Verbatim}[breaklines=true,breakanywhere=true,fontsize=\small]
狗跑去了卫生间。
\end{Verbatim}

由于“狗”具有动物性, “跑”更可能取“快速行走”的含义,
而不是“车辆行驶”的含义.

基于规则方法的优点是解释性强;
缺点是规则编写成本高, 覆盖率有限, 也难以处理开放领域中的复杂语境.

\subsection{基于有监督学习的方法}


有监督 WSD 使用人工标注训练样本.
训练样本中, 目标词在每个上下文中的正确义项已经给出.
模型学习的是:

\begin{Verbatim}[breaklines=true,breakanywhere=true,fontsize=\small]
上下文特征 -> 目标词义项
\end{Verbatim}

课件列出三类代表方法:

\begin{itemize}
\item 基于互信息的方法.
\item 基于贝叶斯分类器的方法.
\item 基于最大熵的方法.
\end{itemize}

有监督方法可以分为两种理解方式:

\begin{itemize}
\item 显式分类器: 直接学习一个分类器, 对每个候选义项打分.
\item 隐式分类指示器: 先学习哪些上下文特征对某个义项有提示作用.
\end{itemize}

\subsubsection{互信息方法}


互信息(mutual information)用来衡量某个上下文指示器与某个义项之间的相关程度.

假设$T_1, T_2, \dots, T_m$是多义词的候选义项,
$V_1, V_2, \dots, V_n$是某个上下文指示器的可能取值.
互信息可理解为:

\[
I(T, V) = \log(P(T, V) / (P(T) P(V)))
\]

如果某个上下文特征$V$与某个义项$T$经常共同出现,
则$I(T,V)$较大, 说明$V$是识别$T$的有效指示器.

使用时, 可以把多个指示器的证据汇总,
选择得到最多支持的义项.

\subsubsection{贝叶斯分类器方法}


贝叶斯方法把词义消歧看作后验概率最大化:

\[
T^* = \arg\max_T P(T | C)
\]

其中$C$表示当前上下文.
常见的朴素贝叶斯近似写作:

\[
P(T | C) prop P(T) \prod_i P(f_i | T)
\]

$f_i$是上下文特征, 如邻近词、主题词、词性、依存关系等.
这种方法实现相对简单, 训练和预测速度快, 但独立性假设较强.

\subsubsection{最大熵方法}


最大熵方法把上下文作为特征信息, 利用最大熵模型对词义进行分类.

在约束条件已知时, 最大熵原则选择不确定性最大、最中立、
没有额外主观假设的概率分布作为模型.
常见形式可以写作:

\[
P(y | x) = (1 / Z(x)) \exp(\sum_i \lambda_i f_i(x, y))
\]

其中:

\begin{itemize}
\item $x$是上下文.
\item $y$是候选义项.
\item $f_i(x,y)$是特征函数.
\item $\lambda_i$是特征权重.
\item $Z(x)$是归一化因子.
\end{itemize}

最大熵的优点是通用性强、无需假设数据分布、能自然融合多种特征;
缺点是训练通常更耗时.

\subsection{基于词典的方法}


基于词典的方法不一定依赖大规模标注语料,
而是利用词典中的定义、义类或翻译信息进行消歧.

课件列出三类:

\begin{table}[H]
\centering
\begingroup
\setlength{\tabcolsep}{2pt}
\small
\caption{基于词典的词义消歧方法}
\begin{tabular}{|p{0.28\textwidth}|p{0.28\textwidth}|p{0.28\textwidth}|}
\hline
方法 & 核心依据 & 直观理解 \\ \hline
基于语义定义 & 词典释义与上下文词重叠 & 哪个义项释义更像当前上下文 \\ \hline
基于义类辞典 & 义项所属语义类与上下文语义类共现 & 上下文更像动物场还是机器场 \\ \hline
基于双语词典 & 翻译后的跨语言共现规律 & 借助另一种语言区分目标义项 \\ \hline
\end{tabular}
\endgroup
\end{table}

基于语义定义的方法常见思路是 Lesk 式重叠:
分析目标词上下文窗口中的词,
再比较这些词与各义项词典注释之间的重合程度.

基于义类辞典的方法利用义项的语义类别.
例如 \texttt{crane} 可表示“鹤”或“起重机”,
分别属于 \texttt{ANIMAL} 和 \texttt{MACHINE}.
若上下文中出现 \texttt{FOOD}、\texttt{EMOTION} 等更常与动物共现的语义类,
则更支持“鹤”的义项;
若出现 \texttt{OPERATE}、\texttt{TOOL} 等语义类, 则更支持“起重机”的义项.

基于双语词典的方法把机器翻译作为桥梁.
例如 \texttt{plant} 可译为“植物”或“工厂”.
在 \texttt{plant life} 中, 可先把 \texttt{life} 译为“生命”,
再统计“生命”与“植物”、“工厂”的共现频率.
若“生命”与“植物”共现更显著, 就判定 \texttt{plant} 为“植物”.
同理, \texttt{manufacturing plant} 中 \texttt{manufacturing} 对应“制造”,
更可能与“工厂”共现, 因此 \texttt{plant} 取“工厂”义项.

\section{语义角色标注}


\subsection{基本概念}


语义角色标注以句子的谓词为中心,
研究句子中各成分与谓词之间的语义关系,
并用语义角色描述这些关系.

可以把它理解为从“句法结构”进一步走向“事件结构”:

\begin{Verbatim}[breaklines=true,breakanywhere=true,fontsize=\small]
谁 对 谁 在 何时 何地 用 什么方式 做了 什么事
\end{Verbatim}

例如:

\begin{Verbatim}[breaklines=true,breakanywhere=true,fontsize=\small]
小明 用 刀 切 苹果
\end{Verbatim}

其中谓词是“切”, 各成分可标注为:

\begin{table}[H]
\centering
\begingroup
\setlength{\tabcolsep}{2pt}
\small
\caption{谓词“切”的语义角色示例}
\begin{tabular}{|p{0.28\textwidth}|p{0.28\textwidth}|p{0.28\textwidth}|}
\hline
成分 & 语义角色 & 说明 \\ \hline
小明 & 施事 / Agent & 动作发出者 \\ \hline
刀 & 工具 / Instrument & 动作使用的工具 \\ \hline
苹果 & 受事 / Patient & 动作影响的对象 \\ \hline
切 & 谓词 / Predicate & 事件核心 \\ \hline
\end{tabular}
\endgroup
\end{table}

\subsection{语义角色与格关系}


语义角色来自语言学中的格语法(case grammar)理论.
这里的“格”指句子中体词与谓词之间的及物性关系,
如动作与施事者、受事者、结果、工具之间的关系.

常见语义角色包括:

\begin{table}[H]
\centering
\begingroup
\setlength{\tabcolsep}{2pt}
\small
\caption{常见语义角色}
\begin{tabular}{|p{0.28\textwidth}|p{0.28\textwidth}|p{0.28\textwidth}|}
\hline
角色 & 含义 & 例子 \\ \hline
施事 Agent & 动作发出者 & 小明切苹果 \\ \hline
受事 Patient & 动作承受者 & 切苹果 \\ \hline
客体 Theme & 被移动、描述或经历的对象 & 移动箱子 \\ \hline
经验者 Experiencer & 感知或心理状态承受者 & 小明喜欢音乐 \\ \hline
工具 Instrument & 动作使用的工具 & 用刀切 \\ \hline
处所 Location & 事件发生地点 & 在厨房做饭 \\ \hline
时间 Time & 事件发生时间 & 昨天到达 \\ \hline
来源 Source & 动作起点 & 从杭州出发 \\ \hline
目标 Goal & 动作终点 & 到上海 \\ \hline
结果 Result & 动作产生的结果 & 打碎杯子 \\ \hline
\end{tabular}
\endgroup
\end{table}

不同语料库或标注体系会使用不同角色集合.
英文领域常见两种方式:

\begin{itemize}
\item 根据动词词典为每个动词规定格框架.
\item 预先定义对多数动词通用的角色集合.
\end{itemize}

中文领域也需要根据语料库规范确定角色标签,
实际使用时要以对应标注体系为准.

\subsection{SRL 的基本流程}


语义角色标注通常包含几个步骤:

\begin{enumerate}
\item[1.] 谓词识别: 找到句子中的谓词.
\item[2.] 谓词义项识别: 判断谓词当前使用的是哪个义项.
\item[3.] 论元识别: 找出哪些成分与谓词有语义关系.
\item[4.] 角色分类: 为每个论元赋予语义角色.
\end{enumerate}

\begin{quote}
\small
\textbf{语义角色标注的一般流程}
\par 节点: srl-input: 输入句子; srl-pred: 谓词识别; srl-arg: 候选论元识别; srl-role: 角色分类; srl-output: 语义角色结构; srl-pred-out: 切、买、给等事件核心; srl-arg-out: 短语或词语候选; srl-role-out: 施事、受事、工具等
\par 边: srl-input -\textgreater{} srl-pred; srl-pred -\textgreater{} srl-arg; srl-arg -\textgreater{} srl-role; srl-role -\textgreater{} srl-output; srl-pred -\textgreater{} srl-pred-out; srl-arg -\textgreater{} srl-arg-out; srl-role -\textgreater{} srl-role-out
\end{quote}

\subsection{方法概述}


课件将 SRL 方法分为三类:

\begin{table}[H]
\centering
\begingroup
\setlength{\tabcolsep}{2pt}
\small
\caption{语义角色标注的三种典型路线}
\begin{tabular}{|p{0.28\textwidth}|p{0.28\textwidth}|p{0.28\textwidth}|}
\hline
方法 & 依赖的上游结构 & 标注对象 \\ \hline
基于短语结构树 & 句法结构分析树 & 短语 \\ \hline
基于依存关系树 & 依存句法树 & 词语 \\ \hline
基于浅层句法语块 & 语块序列 & 浅层语块 \\ \hline
\end{tabular}
\endgroup
\end{table}

三类方法的共同点是:
都需要先借助某种句法或浅层结构来缩小候选论元范围,
再用分类器判断候选成分的语义角色.

\subsection{基于短语结构树的方法}


基于短语结构树的方法利用完全句法分析得到的句法结构树进行 SRL.
其特点是:

\begin{itemize}
\item 标注对象为短语, 如 NP、PP、VP 等.
\item 特征主要从句法结构树中抽取.
\item 适合处理论元以短语形式出现的情况.
\end{itemize}

常用特征包括:

\begin{itemize}
\item 谓词本身及其词性.
\item 候选短语的类型, 如 NP、PP.
\item 候选短语的中心词.
\item 候选短语相对谓词的位置.
\item 句法树中谓词到候选短语的路径.
\item 谓词的语态、子范畴化框架等.
\end{itemize}

常用分类器包括最大熵、SVM、CRF 以及神经网络模型.
早期系统常把“论元识别”和“角色分类”拆成两个分类任务;
现代模型则常把两者联合建模.

\subsection{基于依存关系树的方法}


基于依存关系树的方法利用依存句法分析结果进行 SRL.
其特点是:

\begin{itemize}
\item 标注对象为词语或以词为中心的成分.
\item 特征主要从依存树中抽取.
\item 与“谓词支配哪些词、哪些词依附于谓词”这一语义直觉更接近.
\end{itemize}

课件给出的候选论元搜索过程可以概括为:

\begin{enumerate}
\item[1.] 将谓词作为当前结点.
\item[2.] 将当前结点的所有子结点作为候选论元.
\item[3.] 将当前结点的父结点设为新的当前结点.
\item[4.] 若新当前结点是依存树根结点, 则结束; 否则回到第 2 步.
\end{enumerate}

\begin{quote}
\small
\textbf{基于依存树的候选论元搜索思路}
\par 节点: dep-pred: 谓词结点; dep-child1: 子结点; dep-current: 当前结点; dep-parent: 父结点; dep-cand1: 候选论元; dep-up: 继续向上搜索
\par 边: dep-pred -\textgreater{} dep-child1; dep-pred -\textgreater{} dep-current; dep-current -\textgreater{} dep-parent; dep-child1 -\textgreater{} dep-cand1; dep-parent -\textgreater{} dep-up
\end{quote}

常用特征包括:

\begin{itemize}
\item 谓词和候选词本身.
\item 词性、命名实体类型等词法信息.
\item 候选词与谓词之间的依存关系.
\item 依存路径、路径长度、方向.
\item 候选词在谓词左侧还是右侧.
\item 候选词的父结点、子结点和兄弟结点信息.
\end{itemize}

基于依存树的方法通常比短语结构树方法更直接,
但它高度依赖依存句法分析的准确性.

\subsection{基于浅层句法语块的方法}


基于浅层句法语块的方法只使用语块序列进行 SRL.
它不要求完整句法树, 而是在 Base NP、VP、PP 等语块上进行角色判断.

优点:

\begin{itemize}
\item 上游分析成本低.
\item 对完整句法分析错误的依赖较小.
\item 适合资源有限或需要快速处理的场景.
\end{itemize}

缺点:

\begin{itemize}
\item 缺少深层句法结构信息.
\item 难以处理长距离、嵌套或跨语块的论元关系.
\end{itemize}

可以把三种 SRL 方法理解为同一个问题的三种不同粒度:

\begin{table}[H]
\centering
\begingroup
\setlength{\tabcolsep}{2pt}
\small
\caption{SRL 方法的结构粒度对比}
\begin{tabular}{|p{0.28\textwidth}|p{0.28\textwidth}|p{0.28\textwidth}|}
\hline
粒度 & 上游结构 & 适用特点 \\ \hline
深 & 完整短语结构树 & 结构信息丰富, 分析成本高 \\ \hline
中 & 依存关系树 & 词级关系清晰, 适合事件论元 \\ \hline
浅 & 语块序列 & 速度快, 结构信息较少 \\ \hline
\end{tabular}
\endgroup
\end{table}
